\section{Dataset dan Pengaturan Eksperimen}\label{sec:setup}
\subsection{Dataset}
Eksperimen memakai tiga dataset RCAEval \cite{pham2024rcaeval}, masing-masing mencakup sistem Online Boutique, Sock Shop, dan Train Ticket (Tabel~\ref{tab:datasets}). RE1 hanya berisi metrik, sedangkan RE2 dan RE3 juga berisi log dan trace; penelitian ini hanya memakai channel metrik yang dicantumkan pada Bagian~\ref{sec:method}, sehingga ketiga dataset dapat dibandingkan. RE1 dan RE2 memuat fault sumber daya dan jaringan (RE2 menambahkan fault socket), sedangkan RE3 memuat fault tingkat kode. Setelah pembagian 60/20/20, data uji berisi 75 kasus (RE1), 45 kasus (RE2), dan 18 kasus (RE3). Dataset D1 dari ART tidak dipakai pada versi penelitian ini.
\begin{table}
\caption{Dataset RCAEval yang dipakai dalam penelitian ini \cite{pham2024rcaeval}. Setiap dataset mencakup Online Boutique, Sock Shop, dan Train Ticket.}\label{tab:datasets}
\centering
\small
\begin{tabular}{lrlll}
\toprule
Dataset & Kasus & Kasus per sistem & Telemetri & Tipe fault\\
\midrule
RE1 & 375 & 125 & Hanya metrik & CPU, MEM, DISK, DELAY, LOSS\\
RE2 & 270 & 90 & Metrik, log, trace & Fault RE1 dan SOCKET\\
RE3 & 90 & 30 & Metrik, log, trace & Fault tingkat kode (F1--F5)\\
\bottomrule
\end{tabular}
\end{table}

\todo{Tambahkan jumlah service aplikasi per sistem setelah preprocessing, interval sampling asli tiap dataset, dan jumlah kasus yang kehilangan setiap channel. Konsekuensi: tanpa jumlah service, chance level AC@3 dan Avg@5 tidak dapat dihitung dan AC@3 tidak dapat ditafsirkan.}

\subsection{Rancangan Eksperimen}
Tabel~\ref{tab:config} mencantumkan konfigurasi eksperimen. EX-1 melatih dan menguji pada dataset yang sama. EX-2 melakukan pre-training pada dataset yang berbeda, sehingga setiap target dievaluasi dengan encoder transfer, in-domain, dan acak. Untuk encoder yang dapat dilatih (TS2Vec, TimeMAE, RUN-encoder), setiap konfigurasi disilangkan dengan dua mode fine-tuning dan empat fraksi label, lalu diulang dengan tiga seed pada satu pembagian data tetap.
\begin{table}
\caption{Konfigurasi eksperimen; $\mathcal{S}$ adalah dataset sumber dan $\mathcal{T}$ dataset target. Encoder in-domain pada EX-2 adalah encoder hasil pre-training pada EX-1.}\label{tab:config}
\centering
\small
\begin{tabular}{p{1.9cm}p{3.8cm}p{5.3cm}}
\toprule
 & EX-1 (in-domain) & EX-2 (transfer)\\
\midrule
Pre-training & $\mathcal{S}=\mathcal{T}$ & $\mathcal{S}\neq\mathcal{T}$, $\mathcal{S}\in\{\mathrm{RE1},\mathrm{RE2},\mathrm{RE3}\}$\\
Fine-tuning dan uji & $\mathcal{T}\in\{\mathrm{RE1},\mathrm{RE2},\mathrm{RE3}\}$ & $\mathcal{T}\in\{\mathrm{RE1},\mathrm{RE2}\}$; RE3 bukan target (jenuh)\\
Encoder & Pre-trained, acak & Transfer, in-domain, acak\\
Mode fine-tuning & FineLast, FineAll & FineLast, FineAll\\
Fraksi label & \todo{konfirmasi} & 10\%, 25\%, 50\%, 100\%\\
Seed & 3 & 3\\
Metrik & AC@1 (AC@3, Avg@5) & AC@1 (AC@3, Avg@5); paired bootstrap, 10.000 resampel\\
\bottomrule
\end{tabular}
\end{table}


\subsection{Aturan Perbandingan}\label{sec:fair}
Mekanisme pelatihan ketiga metode berbeda, sehingga satu protokol tunggal tidak dipaksakan kepada semuanya. Eksperimen menjawab tiga pertanyaan, masing-masing dengan aturan perbandingannya sendiri.
\begin{description}
\item[C1: Kinerja RCA (metode terhadap metode).] Setiap metode dijalankan pada konfigurasi terbaiknya yang dipilih pada kasus validasi, dengan kasus, seed, data uji, dan metrik yang sama (Bagian~\ref{sec:metrics}); anggaran tuning tiap metode dilaporkan. RUN-original dan heuristik disertakan sebagai acuan tanpa transfer. Perbandingan ini mengukur apa yang akan diperoleh praktisi dari tiap metode, bukan nilai pre-training.
\item[C2: Efek pre-training dan transfer (di dalam satu metode).] Untuk setiap encoder yang dapat dilatih, encoder hasil pre-training atau transfer dibandingkan dengan encoder acak berarsitektur sama pada head, label, dan mode fine-tuning yang sama, menggunakan paired bootstrap. AC@1 absolut tidak dibandingkan antar metode dalam analisis ini; hanya keuntungan setiap metode terhadap encoder acaknya sendiri yang dibandingkan, dan dinyatakan sebagai efek relatif.
\item[C3: Efek label dan fine-tuning (di dalam satu metode).] AC@1 dilaporkan untuk setiap fraksi label dan mode fine-tuning. RUN-original tidak memakai label dan hanya muncul sebagai garis acuan horizontal.
\end{description}

\subsection{Detail Implementasi}
Tabel~\ref{tab:impl} mencantumkan hiperparameter yang dilaporkan pada paper asli dan nilai yang dipakai di sini. Head, optimizer, aturan early stopping, dan seed dibuat sama pada seluruh kondisi untuk satu metode.
\begin{table}
\caption{Hiperparameter. Baris ``paper'' adalah nilai yang dilaporkan pada paper asli untuk tolok ukurnya masing-masing dan belum tentu sama dengan yang dipakai di sini; nilai ``dipakai'' harus dilengkapi untuk penelitian ini.}\label{tab:impl}
\centering
\scriptsize
\setlength{\tabcolsep}{2.5pt}
\begin{tabular}{p{1.9cm}p{3.0cm}p{3.2cm}p{2.5cm}}
\toprule
Pengaturan & TS2Vec \cite{yue2022ts2vec} & TimeMAE \cite{cheng2026timemae} & RUN \cite{lin2024run}\\
\midrule
Backbone (paper) & Proyeksi masukan, CNN dilated dengan 10 blok residual & Transformer, 8 layer, 4 head, dropout 0,2; decoupled encoder 6 layer & DLinear\\
Ukuran representasi (paper) & 320 (klasifikasi) & 64 & tidak disebutkan\\
Spesifikasi pre-training (paper) & Cropping; timestamp masking dengan $p=0{,}5$ & Rasio mask 60\%; panjang sub-series dicari pada $\{4,8,12,16\}$; ukuran codebook dicari pada $\{64,\dots,512\}$; $\alpha=1$, $\beta$ dicari pada $\{1,\dots,10\}$; momentum 0,99 & Window 32; 50 epoch untuk pre-training dan fine-tuning\\
Optimizer (paper) & Tidak ada di teks utama (appendix) & Adam, lr 0,001, batch 64 & Adam, lr 0,001, batch 128\\
Dipakai pada penelitian ini & \todo{} & \todo{} & \todo{}\\
\midrule
Head, loss, epoch, early stopping & \multicolumn{3}{l}{\todo{disamakan antar metode sejauh encoder memungkinkan}}\\
Perangkat keras, waktu komputasi & \multicolumn{3}{l}{\todo{}}\\
\bottomrule
\end{tabular}
\end{table}

